In this section, we classify the 30 soundness issues according to the affected
feature and the symptom and describe each category with simplified code
snippets.
%
When describing the issues, we analyze the consequences, triggering features,
and status as well.
%
We identify five affected features and ten symptoms, as shown in
\cref{tab:features};
%
we identify four consequences, as shown in \cref{tab:consequences}, and fourteen
triggering features:
%
{\bf implied} bounds, {\bf assoc}iated types, \C{impl} traits ({\bf IT}s), {\bf
RPIT}s, {\bf tr}ait {\bf obj}ects, higher-rank ({\bf HR}) polymorphism,
associated {\bf const}ants, function pointers ({\bf fn ptr}s), {\bf clo}sures,
{\bf rec}ursive calls, {\bf async}hrony, {\bf try} propagation, {\bf arr}ays,
and {\bf enum}s
%
(in tables, we abbreviate them as shown in bold).
%
Note that these affected/triggering features are not intended to systematically
enumerate all features in Rust;
%
rather, they aid in understanding the soundness issues we collected.
%
We also classify status into four kinds, as shown in \cref{tab:status-kind}.
%
Note that an issue is considered fixed if the code given in the issue is
rejected by the latest stable compiler, regardless of its GitHub issue
resolution status.
%
We use this criterion because an issue may reveal multiple compiler bugs and may
not be closed even after the soundness bug has been fixed due to other remaining
bugs.
%
\cref{tab:issues} summarizes the characteristics of each issue.
%
For the sake of brevity, we only present code snippets for some issues;
%
the omitted ones are available in our supplementary material~\cite{supp}.

\begin{table}[t]
\caption{Classification of affected features and symptoms of soundness issues}
\label{tab:features}
\vspace{-1.15em}
\renewcommand{\arraystretch}{0.9}
  {\scriptsize
\begin{tabular}{l|l}
  Affected feature & Symptoms \\ \hline
  Lifetimes & Missing lifetime checks, Anonymous lifetimes allowed \\
  Traits & Overlapping impls, Missing dyn-compat checks, Orphan impls, Incorrect trait solving \\
  Type inference & Arbitrary type argument selection, Arbitrary hidden type selection \\
  Type well-formedness & Missing well-formedness checks \\
  Type cast & Incorrect type cast allowed \\
\end{tabular}
  }
\vspace{-0.65em}
\end{table}

\begin{table}[t]
\caption{Classification of consequences of soundness issues}
\label{tab:consequences}
\vspace{-1.15em}
\renewcommand{\arraystretch}{0.9}
  {\scriptsize
\begin{tabular}{l|l}
  Consequence & Description \\ \hline
  Memory bug & A memory bug can occur at run time \\
  Breaking change & A benign change in a library can lead to a compilation error in a client code \\
  Inconsistency & A code snippet passes type checking while similar ones are rejected \\
  Misleading behavior & The run-time behavior could be misleading \\
\end{tabular}
  }
\vspace{-0.65em}
\end{table}

\begin{table}[t]
\caption{Classification of status of soundness issues}
\label{tab:status-kind}
\vspace{-1.15em}
\renewcommand{\arraystretch}{0.9}
  {\scriptsize
\begin{tabular}{l|l}
  Status & Description \\ \hline
  Fixed & The compiler has been fixed, and the code snippet is rejected by the latest release (1.93.0) \\
  Confirmed & The C-bug or I-unsound label has been assigned to by a person other than the reporter\\
  Reported & The issue is open but not fixed or confirmed yet \\
  Non-bug & The issue has been closed without being fixed or confirmed \\
\end{tabular}
  }
\vspace{-1.0em}
\end{table}

\begin{table}[t]
\caption{Soundness issues in \rustc{}}
\label{tab:issues}
  \renewcommand{\arraystretch}{0.9}
  \vspace{-1.0em}
  {  \scriptsize
\begin{tabular}{l|l|l|l|l|l}
Issue \# & Affected Feature & Symptom & Consequence & Triggers & Status \\ \hline
\href{https://github.com/rust-lang/rust/issues/96460}{96460}$^\mathsection$& Type inference & Hidden type & Inconsistency & RPIT & Fixed \\
\href{https://github.com/rust-lang/rust/issues/96927}{96927}&Type inference&Type argument&Inconsistency&try&Confirmed$^\mathparagraph$\\
\href{https://github.com/rust-lang/rust/issues/98117}{98117}$^\mathsection$& Lifetimes & Missing lifetime checks & Memory bug & implied, tr obj & Fixed \\
\href{https://github.com/rust-lang/rust/issues/98543}{98543}$^\ddag$&Lifetimes&Missing lifetime checks&Memory bug&implied, assoc&Fixed\\
\href{https://github.com/rust-lang/rust/issues/99158}{99158}$^\ast$&Lifetimes&Missing lifetime checks&Memory bug&implied, fn ptr&Confirmed\\
\href{https://github.com/rust-lang/rust/issues/99554}{99554}&Traits&Orphan&Breaking
change&assoc&Fixed$^{\ast\ast}$\\
\href{https://github.com/rust-lang/rust/issues/100051}{100051}&Lifetimes&Missing lifetime checks&Memory bug&implied, assoc, HR&Reported\\
\href{https://github.com/rust-lang/rust/issues/104005}{104005}$^\mathsection$& Lifetimes & Missing lifetime checks & Memory bug & implied, tr obj & Fixed \\
\href{https://github.com/rust-lang/rust/issues/104763}{104763}$^\mathsection$& Lifetimes & Missing lifetime checks & Inconsistency & const & Fixed \\
\href{https://github.com/rust-lang/rust/issues/104764}{104764}$^\mathsection$& Lifetimes & Missing lifetime checks & Inconsistency & assoc & Fixed \\
\href{https://github.com/rust-lang/rust/issues/105251}{105251}$^\ddag$&Lifetimes&Missing lifetime checks&Inconsistency&assoc, RPIT&Fixed\\
\href{https://github.com/rust-lang/rust/issues/105787}{105787}$^\ddag$&Traits&Overlapping&Inconsistency&assoc, HR&Fixed\\
\href{https://github.com/rust-lang/rust/issues/106040}{106040}&Traits&Incorrect solving&Inconsistency&assoc&Reported\\
\href{https://github.com/rust-lang/rust/issues/107268}{107268}&Type inference&Type argument&Inconsistency&assoc, fn ptr&Fixed\\
\href{https://github.com/rust-lang/rust/issues/108468}{108468}&Lifetimes&Anonymous&Inconsistency&IT, async&Reported\\
\href{https://github.com/rust-lang/rust/issues/109628}{109628}&Lifetimes&Missing lifetime checks&Inconsistency&assoc&Fixed\\
\href{https://github.com/rust-lang/rust/issues/110925}{110925}&Type wf-ness&Missing wf-ness checks&Inconsistency&RPIT&Reported\\
\href{https://github.com/rust-lang/rust/issues/111935}{111935}$^\ddag$&Type inference&Hidden type&Inconsistency&RPIT, rec&Fixed\\
\href{https://github.com/rust-lang/rust/issues/112417}{112417}&Lifetimes&Missing lifetime checks&Memory bug&implied, RPIT&Fixed\\
\href{https://github.com/rust-lang/rust/issues/114061}{114061}&Traits&Overlapping&Inconsistency&assoc, HR&Fixed\\
\href{https://github.com/rust-lang/rust/issues/114389}{114389}$^\dag$&Traits&Overlapping&Memory bug&assoc, tr obj&Confirmed\\
\href{https://github.com/rust-lang/rust/issues/118876}{118876}$^\mathsection$& Lifetimes & Missing lifetime checks & Memory bug & implied, tr obj, HR, clo & Fixed \\
\href{https://github.com/rust-lang/rust/issues/123312}{123312}&Lifetimes&Missing lifetime checks&Inconsistency&clo&Reported\\
\href{https://github.com/rust-lang/rust/issues/126079}{126079}$^\mathsection$& Traits & Missing dyn-compat checks & Memory bug & assoc, tr obj & Fixed \\
\href{https://github.com/rust-lang/rust/issues/133361}{133361}&Traits&Overlapping&Memory bug&assoc, tr obj&Fixed\\
\href{https://github.com/rust-lang/rust/issues/134311}{134311}&Type inference&Hidden type&Inconsistency&assoc, RPIT&Reported\\
\href{https://github.com/rust-lang/rust/issues/136508}{136508}&Type cast&Incorrect cast&Misleading&enum&Reported\\
\href{https://github.com/rust-lang/rust/issues/139406}{139406}&Type inference&Hidden type&Inconsistency&RPIT, rec&Fixed\\
\href{https://github.com/rust-lang/rust/issues/141713}{141713}&Lifetimes&Missing lifetime checks&Memory bug&assoc, HR, clo&Confirmed\\
\href{https://github.com/rust-lang/rust/issues/144045}{144045}&Lifetimes&Missing lifetime checks&Inconsistency&arr&Reported\\
\end{tabular}
  \\[-0.1em] $^\ast$ A duplicate of Issue 25860, which was reported in 2015.
  \\[-0.1em] $^\dag$ A duplicate of Issue 57893, which was reported in 2019.
  \\[-0.1em] $^\ddag$ Overlaps between ours and Liu et al.'s.
  \\[-0.1em] $^\mathsection$ Added after comparison with Liu et al.'s.
  \\[-0.1em] $^\mathparagraph$ This issue was closed as a duplicate of Issue 94902,
  which is not itself a soundness issue but shares the same root cause.
  Since Issue 94902 satisfies the ``confirmed'' criterion, Issue 96927 is
  classified as confirmed.
  \\[-0.6em] $^{\ast\ast}$ A warning occurs instead of an error.
  }
\vspace{-1.5em}
\end{table}

\subsection{Lifetimes}

\subsubsection{Missing Lifetime Checks}

In fourteen soundness issues, the compiler does not check lifetimes properly.
%
As explained in \cref{sec:background:lifetime}, a reference can be used wherever
its lifetime outlives the expected lifetime, and the compiler must check
lifetimes when references are used to enforce this requirement.
%
If the compiler fails to check lifetimes, references may be used even after
their lifetimes have ended, incurring memory bugs like use-after-free.
%
Among the fourteen issues, eight issues actually break memory safety, while six
issues do not result in such misuse of references despite missing lifetime
checks, only exhibiting inconsistent type checking results.
%
We first discuss the memory-bug-incurring issues and then the others.

{\it Issue 99158.}
%
In this issue, the compiler does not check implied bounds when coercing function
definitions to function pointers.
%
In \cref{lst:i99158}, \C{foo} has the implied bound \C{'a: 'static} due to its
second parameter (line~\ref{line:99158:1}).
%
However, when \C{bar} assigns \C{foo} to \C{f}, this implied bound is not
checked, and the assignment succeeds even though \C{bar} does not have \C{'a:
'static} (line~\ref{line:99158:3}).
%
By calling \C{f}, \C{bar} can assign \C{'static} to any reference;
%
if this reference is not valid for the entire execution, dereferencing it would
lead to a memory bug.
%
Note that explicitly adding the bound \C{'a: 'static} to \C{foo} causes the
assignment to be rejected correctly.

In six other issues, the compiler does not check implied bounds in other
contexts and fails to prevent memory bugs:
%
it does not check implied bounds when
%
a trait has a lifetime bound (Issue 98117),
%
a \emph{self type} (i.e., \C{S} in \C{<S as X>::A}) introduces implied bounds
(Issues 98543 and 100051),
%
an implementing type in an \C{impl} introduces implied bounds (Issue 104005),
%
an RPIT hidden type is determined (Issue 112417),
%
or a closure is coerced (Issue 118876).

{\it Issue 141713.}
%
In this issue, the compiler does not properly check the lifetime of the associated
type of an HRTB.
%
In \cref{lst:i141713}, \C{X<T>} declares \C{Fn() -> T} as a \emph{supertrait},
i.e., a type implementing \C{X<T>} must also implement \C{Fn() -> T}
(line~\ref{line:i141713:1}).
%
\C{Fn() -> T} is a built-in trait automatically implemented by every closure
that takes no parameters and returns \C{T}.
%
It has an associated type named \C{Output}, which is set to \C{T}.
%
Additionally, every type implementing \C{Fn() -> T} also implements \C{X<T>}
(line~\ref{line:i141713:2}).
%
In \C{foo}, \C{F} has an HRTB, but it is unclear what \C{F::Output} denotes;
%
the compiler allows using \C{F::Output} as a reference with \C{'static} in
\C{foo} (line~\ref{line:i141713:3}), while it allows using a reference with any
lifetime as \C{F::Output} in \C{bar} (line~\ref{line:i141713:4}).
%
Note that \C{|| \&0} is a closure that takes no parameters and returns \C{\&0}.

{\it Issue 144045.}
%
In this issue, the compiler does not check the lifetime of a length-0 array.
%
In \cref{lst:i144045}, \C{foo} returns \C{[v; 1]}, which is a length-1 array
with \C{v} as its element, and it is rejected because \C{'a} does not outlive
\C{'static} (line~\ref{line:144045:1}).
%
On the other hand, \C{bar} is accepted, even though the only difference from
\C{foo} is that it returns \C{[v; 0]}, which is a length-0 array
(line~\ref{line:144045:2}).
%
This issue does not cause memory bugs because \C{bar} does not actually return
any reference.
%
Nevertheless, this behavior might seem surprising, as developers usually expect
the type checking of array types to be independent of their lengths.

Similarly, five other issues also reveal cases where the compiler does not check
lifetimes but still prevents memory bugs.
%
In particular, it does not check lifetimes when
%
an associated constant or type is accessed (Issues 104763 and 104764),
%
an RPIT has nested \C{impl} traits (Issue 105251),
%
a type parameter's associated type has a trait bound (Issue 109628),
%
or a struct value is constructed inside a closure (Issue 123312).

\begin{figure}[t]
\begin{lstlisting}[caption={Issue 99158}, label={lst:i99158}]
fn foo<'a>(v: &'a i32, u: &'static &'a i32) -> &'static i32 { v } $\label{line:99158:1}$
fn bar<'a>(v: &'a i32) -> &'static i32 { $\label{line:99158:2}$
    let f: fn(&'a i32, _) -> &'static i32 = foo; f(v, &&0) } $\label{line:99158:3}$
\end{lstlisting}
\vspace{-1.2em}
\begin{lstlisting}[caption={Issue 141713}, label={lst:i141713}]
trait X<T>: Fn() -> T {} $\label{line:i141713:1}$
impl<T, F: Fn() -> T> X<T> for F {} $\label{line:i141713:2}$
fn foo<F>(v: F, u: F::Output) -> &'static i32 where for<'a> F: X<&'a i32>{u}$\label{line:i141713:3}$
fn bar<'a>(v: &'a i32) -> &'static i32 { foo(|| &0, v) } $\label{line:i141713:4}$
\end{lstlisting}
\vspace{-1.2em}
\begin{lstlisting}[caption={Issue 144045}, label={lst:i144045}]
fn foo(v: &'a i32) -> [&'static i32; 1] { [v; 1] } // error $\label{line:144045:1}$
fn bar(v: &'a i32) -> [&'static i32; 0] { [v; 0] } // ok $\label{line:144045:2}$
\end{lstlisting}
\vspace{-1.2em}
\begin{lstlisting}[caption={Issue 108468}, label={lst:i108468}]
trait X<T> {} $\label{line:i108468:1}$
fn foo(v: impl X<&i32>) {} // error $\label{line:i108468:2}$
fn bar<'a>(v: impl X<&'a i32>) {} // ok $\label{line:i108468:3}$
async fn baz(v: impl X<&i32>) {} // ok $\label{line:i108468:4}$
\end{lstlisting}
\vspace{-1.2em}
\end{figure}

\subsubsection{Anonymous Lifetimes Allowed}

In one soundness issue, the compiler unexpectedly allows \emph{anonymous
lifetimes}.
%
In Rust, when the compiler can automatically introduce necessary lifetime
parameters according to predefined rules~\cite{rust-lifetime-elision},
programmers can use anonymous lifetimes.
%
For example, \C{foo} and \C{bar} are equivalent in the following code:

\vspace{-0.5em}
\begin{lstlisting}
fn foo(v: &i32) {} fn bar<'a>(v: &'a i32) {}
\end{lstlisting}
\vspace{-0.5em}

\noindent
%
However, anonymous lifetimes are explicitly disallowed in some contexts, and if
the compiler allows them, it may be considered a soundness bug.

{\it Issue 108468.}
%
In this issue, using an \C{async} function causes the compiler to allow
anonymous lifetimes.
%
In \cref{lst:i108468}, the function \C{foo} is rejected because anonymous
lifetimes must not occur in \C{impl} traits (line~\ref{line:i108468:2}), while
\C{bar} is accepted (line~\ref{line:i108468:3}).
%
However, the \C{async} function \C{baz} passes type checking even though it
contains an anonymous lifetime in an \C{impl} trait (line~\ref{line:i108468:4}).
%
% In fact, when the unstable feature \C{anonymous\_lifetime\_in\_impl\_trait} is
% explicitly enabled, such anonymous lifetimes are allowed even for synchronous
% functions.
%
% In the future, the \rustc{} developers may stabilize this feature, making
% \C{foo} accepted as well, rather than fixing the compiler to reject \C{baz}.

\subsection{Traits}

\subsubsection{Overlapping Implementations}

In four soundness issues, the compiler allows \emph{overlapping impls}, i.e.,
those implementing the same trait for the same type.
%
Rust prohibits overlapping \C{impl}s because they create ambiguity when deciding
which \C{impl} to use during method resolution.
%
For example, the following code is rejected:

\vspace{-0.5em}
\begin{lstlisting}
trait X { fn foo(self) -> i32; } struct S;
impl X for S{fn foo(self)->i32{0}}impl X for S{fn foo(self)->i32{1}} // error
\end{lstlisting}
\vspace{-0.5em}

\noindent
%
Therefore, if code with overlapping \C{impl}s passes type checking, it can be
considered a soundness bug.
%
Among the four issues in this category, two lead to memory bugs, while two
only cause inconsistency.

{\it Issue 114389.}
%
In this issue, the compiler allows a user-defined \C{impl} to overlap with a
\emph{built-in impl}---for every trait \C{X}, the type \C{dyn X} automatically
implements \C{X}, called a built-in \C{impl}.
%
In \cref{lst:i114389}, every type implements \C{X}, even including dynamically
sized types due to the bound \C{?Sized} on \C{T} (line~\ref{line:114389:2}).
%
In \C{foo}, since \C{<T as X>::A} and \C{\&'static i32} are the same type, it
can return the given value as is (line~\ref{line:114389:3}).
%
On the other hand, when \C{bar} calls \C{foo}, the type argument is \C{dyn X<A =
\&'a i32>}, so the parameter type \C{<T as X>::A} becomes \C{\&'a i32}.
%
Consequently, \C{bar} can use a value of \C{\&'a i32} as \C{\&'static i32} by
passing it to \C{foo} (line~\ref{line:114389:4}).
%
This issue arises because \C{dyn X<A = \&'a i32>} has two overlapping \C{impl}s
of \C{X}:
%
the one with \C{A = \&'static i32} defined in line~\ref{line:114389:2} and the
built-in \C{impl} with \C{A = \&'a i32}.

In another issue (Issue 133361), two user-defined \C{impl}s for trait objects
overlap and similarly incur memory bugs by setting the same associated type to
two different types.
%
In two other issues, the compiler fails to detect overlapping \C{impl}s when
higher-rank polymorphism is involved, but memory bugs do not occur
%
because the compiler always selects one \C{impl} and ignores the other (Issue
105787),
%
or when the associated type is used, the compiler recognizes that two different
types are possible and rejects the code (Issue 114061).

\begin{figure}[t]
\begin{lstlisting}[caption={Issue 114389}, label={lst:i114389}]
trait X { type A; } impl<T: ?Sized> X for T { type A = &'static i32; } $\label{line:114389:2}$
fn foo<T: ?Sized>(v: <T as X>::A) -> &'static i32 { v } $\label{line:114389:3}$
fn bar<'a>(v: &'a i32) -> &'static i32 { foo::<dyn X<A = &'a i32>>(v) } $\label{line:114389:4}$
\end{lstlisting}
\vspace{-2em}
\end{figure}

\subsubsection{Missing Dyn-compatibility Checks}

In one soundness issue (Issue 126079), the compiler does not check the
\emph{dyn-compatibility} of a trait.
%
To use a type \C{dyn T}, the trait \C{T} must satisfy dyn-compatibility, which
requires several conditions, such as not having a method return type containing
\C{Self}.
%
However, the compiler does not check dyn-compatibility when an associated type
of a supertrait, referred to using \C{Self}, is used as a return type.
%
This allows a single return type to be interpreted as two different types,
leading to memory bugs.

\subsubsection{Orphan Implementations}

In one soundness issue, the compiler allows \emph{orphan impls}, i.e., those
implementing external traits for external types.
%
In Rust, a \emph{crate} is a compilation unit, and items defined in other crates
are considered external.
%
If orphan \C{impl}s are allowed, a seemingly innocent change in a library can
prevent a dependent program from compiling:

\vspace{-0.5em}
\begin{lstlisting}
// crate a $\label{line:orphan:1}$
trait X {} struct S; $\label{line:orphan:2}$
// crate b (depends on crate a) $\label{line:orphan:3}$
trait Y {} struct R; $\label{line:orphan:4}$
impl a::X for a::S {} // error $\label{line:orphan:5}$
impl a::X for R {} impl Y for a::S {} // ok $\label{line:orphan:6}$
\end{lstlisting}
\vspace{-0.5em}

\noindent
%
While the crate \C{a} currently does not implement \C{X} for \C{S}, it can
freely add an \C{impl} in the future.
%
If the crate \C{b} is allowed to implement \C{X} for \C{S}, such an update in
\C{a} becomes an abrupt breaking change, making \C{b} uncompilable due to two
overlapping \C{impl}s of \C{X} for \C{S}.
%
To prevent this in advance, the orphan implementation is rejected
(line~\ref{line:orphan:5}), while implementing \C{X} for its own type \C{R} or
its own trait \C{Y} for \C{S} is accepted (line~\ref{line:orphan:6}).
%
Therefore, the issue in this category leads to a failure in preventing future
breaking changes.

{\it Issue 99554.}
%
In this issue, the compiler allows an orphan \C{impl} when an associated type is
used.
%
In \cref{lst:i99554}, \C{a} defines \C{X} (line~\ref{line:i99554:2}).
%
In \C{b}, implementing \C{X} for any type \C{T} that implements \C{Y} is
considered orphan and correctly rejected, as such \C{T} can be external
(line~\ref{line:i99554:5}).
%
However, writing \C{<T as Y>::A}, which is equivalent to \C{T}, causes the
compiler to accept the \C{impl} (line~\ref{line:i99554:6}).

\subsubsection{Incorrect Trait Solving}

In one soundness issue, the compiler draws an incorrect conclusion in
\emph{trait solving}, i.e., deciding whether a given type implements a certain
trait.

{\it Issue 106040.}
%
In this issue, trait solving succeeds despite the existence of cyclic
requirements when associated types have trait bounds.
%
In \cref{lst:i106040}, the compiler needs to determine whether \C{S} implements
\C{Unpin} (line~\ref{line:i106040:4}), where \C{Unpin} is a trait defined in the
standard library.
%
This requires \C{S: X<A = S>} (line~\ref{line:i106040:3}), which in turn
requires \C{S: Unpin} (line~\ref{line:i106040:2}), forming a cycle.
%
However, the compiler concludes that \C{S} implements \C{Unpin} and accepts the
entire program.
%
This result can be considered inconsistent because the compiler successfully
detects the cycle and rejects the code when the bound is directly applied to the
type parameter, instead of the associated type, as follows:

\vspace{-0.5em}
\begin{lstlisting}[firstnumber=2]
impl<T: Unpin> X for T { type A = T; }
\end{lstlisting}
\vspace{-0.5em}

\begin{figure}[t]
\begin{lstlisting}[caption={Issue 99554}, label={lst:i99554}]
// crate a $\label{line:i99554:1}$
trait X<T, U> {} $\label{line:i99554:2}$
// crate b $\label{line:i99554:3}$
struct S; trait Y { type A; } impl<T> Y for T { type A = T; } $\label{line:i99554:4}$
impl<T: Y> a::X<S, T> for T {} // error $\label{line:i99554:5}$
impl<T: Y> a::X<S, T> for <T as Y>::A {} // ok $\label{line:i99554:6}$
\end{lstlisting}
\vspace{-1.2em}
\begin{lstlisting}[caption={Issue 106040}, label={lst:i106040}]
trait X { type A where Self: Unpin; } $\label{line:i106040:1}$
impl<T> X for T { type A = T where T: Unpin; } $\label{line:i106040:2}$
struct S; impl Unpin for S where S: X<A = S> {} $\label{line:i106040:3}$
fn foo<T: Unpin>() {} fn bar() { foo::<S>() } $\label{line:i106040:4}$
\end{lstlisting}
\vspace{-1.2em}
\begin{lstlisting}[caption={Issue 96927}, label={lst:i96927}]
fn foo<T>() -> Option<T> { None } $\label{line:i96927:1}$
fn bar() -> Option<i32> { foo(); None } // error $\label{line:i96927:2}$
fn baz() -> Option<i32> { foo()?; None } // ok $\label{line:i96927:3}$
\end{lstlisting}
\vspace{-2em}
\end{figure}

\subsection{Type Inference}

\subsubsection{Arbitrary Type Argument Selection}

In two soundness issues, the compiler arbitrarily chooses a type argument during
type inference even when multiple types are possible.
%
In Rust, type inference of a type argument should succeed only when the possible
type is unique.
%
In the following example, \C{foo(0)} is accepted because \C{i32} is the only
possible type, while \C{baz()} is rejected because any type would be valid:

\vspace{-0.5em}
\begin{lstlisting}
fn foo<T>(t: T) {} fn bar() { foo(0); } // ok
fn baz<T>() {} fn qux() { baz(); } // error
\end{lstlisting}
\vspace{-0.5em}

\noindent
%
If the compiler arbitrarily chooses one among multiple possible types, it can be
considered inconsistent, and the run-time behavior may deviate from the
developer's expectation.

{\it Issue 96927.}
%
In this issue, type inference selects an arbitrary type when the try propagation
operator (\C{?}) is used.
%
In \cref{lst:i96927}, \C{foo} returns \C{Option<T>}, which expresses the
optional existence of a value, allowing \C{Some(v)} for a value or \C{None} for
no value (line~\ref{line:i96927:1}).
%
Since \C{bar} does not use the return value, the type argument can be any type,
and type inference fails (line~\ref{line:i96927:2}).
%
On the other hand, \C{baz} uses \C{?}, which immediately returns \C{None} if the
result is \C{None} and extracts the inner value otherwise
(line~\ref{line:i96927:3}).
%
The compiler still has no information about the type, but chooses an arbitrary
one---\C{!} (signifying that the function never returns) in recent versions, or
\C{()} (the empty tuple) in older versions---and accepts the code.

In another issue (Issue 107268), the compiler chooses an arbitrary type as a
type argument when a function with a lifetime parameter is coerced to a function
pointer.

\subsubsection{Arbitrary Hidden Type Selection}

In four soundness issues, type inference chooses an arbitrary type as the
hidden type for an RPIT even when multiple types are possible.
%
Such issues are also considered inconsistent and may lead to unexpected run-time
behavior.

{\it Issue 139406.}
%
In this issue, type inference selects an arbitrary hidden type when recursive
calls exist.
%
The function \C{foo} passes the result of a recursive call to another recursive
call (line~\ref{line:i139406:3}).
%
While both \C{T} and \C{!} are valid choices, the compiler chooses \C{T} and
accepts the code.

In three other issues, the hidden type is arbitrarily selected when
a built-in type \C{PhantomData} (Issue 96460),
recursive calls (Issue 111935), or associated types (Issue 134311) are involved.

\begin{figure}[t]
\begin{lstlisting}[caption={Issue 139406}, label={lst:i139406}]
trait X {} impl<T> X for T {} $\label{line:i139406:1}$
fn foo<T>(v: T) -> impl X { foo::<T>(foo::<T>(v)) } $\label{line:i139406:3}$
\end{lstlisting}
\vspace{-1em}
\begin{lstlisting}[caption={Issue 110925}, label={lst:i110925}]
trait X<T> {} impl<T> X<T> for i32 {} trait Y {} struct S<T: Y>(T); $\label{line:i110925:1}$
fn foo() -> impl X<S<i32>> { 0 } // ok $\label{line:i110925:2}$
fn bar(v: impl X<S<i32>>) {} // error $\label{line:i110925:3}$
\end{lstlisting}
\vspace{-1em}
\begin{lstlisting}[caption={Issue 136508}, label={lst:i136508}]
enum E { A } enum F { B(i32) } $\label{line:i136508:1}$
fn foo() { let x = E::A as i32; // ok $\label{line:i136508:2}$
           let y = F::B(0) as i32; // error $\label{line:i136508:3}$
           let z = F::B as i32; } // ok $\label{line:i136508:4}$
\end{lstlisting}
\vspace{-1em}
\end{figure}

\subsection{Type Well-formedness}

{\it Issue 110925.}
%
In this issue, the compiler does not check the well-formedness of a type in an
RPIT.
%
In \cref{lst:i110925}, \C{S<i32>} is not well-formed because no type implements
\C{Y}, while \C{S<T>} requires \C{T: Y} (line~\ref{line:i110925:1}).
%
However, \C{foo} is accepted despite using \C{S<i32>} in the RPIT
(line~\ref{line:i110925:2}).
%
The compiler rejects the use of \C{S<i32>} in other contexts, including APITs
(line~\ref{line:i110925:3}).
%
This inconsistency does not lead to memory bugs because a value of \C{S<i32>}
cannot be constructed.

\subsection{Type Cast}

{\it Issue 136508.}
%
This issue allows casting an enum constructor to an integer.
%
In Rust, an enum consists of multiple \emph{variants}, each of which may have
associated data.
%
In \cref{lst:i136508}, \C{E} has one variant \C{A} without data, and \C{F} has
one variant \C{B} with an \C{i32} value (line~\ref{line:i136508:1}).
%
Rust allows casting \C{E::A}, which has no data, to an integer
(line~\ref{line:i136508:2}), but disallows casting \C{F::B(0)}, which contains
data (line~\ref{line:i136508:3}).
%
Somewhat surprisingly, casting \C{F::B} to an integer is accepted
(line~\ref{line:i136508:4}).
%
In fact, this behavior is intended according to the Rust Reference, as \C{F::B}
is an enum constructor and treated as a function;
%
casting it to an integer gives its address.
%
However, the reporter of this issue claims that this behavior is misleading
because a developer might attempt to cast \C{F::B} under the mistaken impression
that it is a variant without data.
